
\documentclass{article}
\usepackage{colortbl}
\usepackage{makecell}
\usepackage{multirow}
\usepackage{supertabular}

\begin{document}

\newcounter{utterance}

\centering \large Interaction Transcript for game `dond', experiment `semi', episode 2 with Qwen3.6{-}35B{-}A3B{-}FP8{-}without{-}reasoning.
\vspace{24pt}

{ \footnotesize  \setcounter{utterance}{1}
\setlength{\tabcolsep}{0pt}
\begin{supertabular}{c@{$\;$}|p{.15\linewidth}@{}p{.15\linewidth}p{.15\linewidth}p{.15\linewidth}p{.15\linewidth}p{.15\linewidth}}
    \# & $\;$A & \multicolumn{4}{c}{Game Master} & $\;\:$B\\
    \hline

    \theutterance \stepcounter{utterance}  
    & & \multicolumn{4}{p{0.6\linewidth}}{
        \cellcolor[rgb]{0.9,0.9,0.9}{
            \makecell[{{p{\linewidth}}}]{
                \texttt{\tiny{[P1$\langle$GM]}}
                \texttt{Igrate igro sodelovalnega pogajanja, v kateri se morate dogovoriti s drugim igralcem o tem, kako razdeliti nabor predmetov med seboj.} \\
\\ 
\texttt{Pravila:} \\
\texttt{(a) Ti in drugi igralec dobite nabor predmetov. Vsakemu od vas je dana tudi tajna funkcija vrednosti, ki predstavlja, kako zelo cenite vsako vrsto predmeta.} \\
\texttt{(b) Izmenjujete si sporočila z drugim igralcem, da se dogovorite, kdo dobi katere predmete. Vsak lahko pošlje največ 5 sporočil ali zaključi igro prej tako, da v katerem koli trenutku daš skritega predloga.} \\
\texttt{(c) Prosimo vas, da pošljete skritega predloga, ki navaja predmete, ki jih želite, oblikovan v oglatih oglatih oklepajih, kot sledi: "{[}Predlog: <število> <ime predmeta>, <število> <ime predmeta>, <...>{]}"} \\
\texttt{(d) Če so vaši predlogi dopolnilni, tj. je dovolj predmetov za izpolnitev obeh predlogov, vsak igralec dobi rezultat na podlagi vsote vrednosti predmetov, ki jih je prejel. V nasprotnem primeru oba dobita nič točk.} \\
\texttt{(e) Vaš cilj je povečati vsoto vašega rezultata in rezultata drugega igralca.} \\
\\ 
\texttt{Začnimo.} \\
\\ 
\texttt{Niz razpoložljivih predmetov je:} \\
\\ 
\texttt{1 torba, 3 luči, 1 cvet} \\
\\ 
\texttt{Vaša tajna funkcija vrednosti je:} \\
\\ 
\texttt{torba: 3, luč: 0, cvet: 7} \\
\\ 
\texttt{POMEMBNO: Vaša sporočila, razen če gre za skriti predlog, se posredujejo neposredno drugemu igralcu, zato ne vključujte nobenega odgovora na pravila ali besedilo, ki napoveduje vaše sporočilo. Če želite podati skriti predlog, uporabite navedeni format. Ne uporabljajte oglatih oglatih oklepajev pri komunikaciji z drugim igralcem, saj bo to razumljeno kot vaš skritega predloga.} \\
            }
        }
    }
    & \\ \\

    \theutterance \stepcounter{utterance}  
    & \multicolumn{4}{p{0.6\linewidth}}{
        \cellcolor[rgb]{0.9,0.9,0.9}{
            \makecell[{{p{\linewidth}}}]{
                \texttt{\tiny{[P1$\rangle$GM]}}
                \texttt{Pozdravljen! Predlagam, da se dogovorimo za razdelitev, ki bo čim bolj izkoristila vrednost predmetov za oba nas. Glede na to, da ti torba in cvet verjetno veliko pomenita, predlagam naslednjo porazdelitev:} \\
\\ 
\texttt{Ti dobiš 1 torbo in 1 cvet.} \\
\texttt{Jaz dobim 3 luči.} \\
\\ 
\texttt{To je po mojem mnenju pravična razdelitev, saj vsak dobi predmete, ki jih ima v svojem naboru, in se izognemo spopadom. Ali se strinjaš s tem predlogom?} \\
            }
        }
    }
    & & \\ \\

    \theutterance \stepcounter{utterance}  
    & & & \multicolumn{2}{p{0.3\linewidth}}{
        \cellcolor[rgb]{0.9,0.9,0.9}{
            \makecell[{{p{\linewidth}}}]{
                \texttt{\tiny{[GM$|$GM]}}
                \texttt{continue} \\
            }
        }
    }
    & & \\ \\

    \theutterance \stepcounter{utterance}  
    & & & \multicolumn{2}{p{0.3\linewidth}}{
        \cellcolor[rgb]{0.9,0.9,0.9}{
            \makecell[{{p{\linewidth}}}]{
                \texttt{\tiny{[GM$|$GM]}}
                \texttt{Pozdravljen! Predlagam, da se dogovorimo za razdelitev, ki bo čim bolj izkoristila vrednost predmetov za oba nas. Glede na to, da ti torba in cvet verjetno veliko pomenita, predlagam naslednjo porazdelitev:} \\
\\ 
\texttt{Ti dobiš 1 torbo in 1 cvet.} \\
\texttt{Jaz dobim 3 luči.} \\
\\ 
\texttt{To je po mojem mnenju pravična razdelitev, saj vsak dobi predmete, ki jih ima v svojem naboru, in se izognemo spopadom. Ali se strinjaš s tem predlogom?} \\
            }
        }
    }
    & & \\ \\

    \theutterance \stepcounter{utterance}  
    & & \multicolumn{4}{p{0.6\linewidth}}{
        \cellcolor[rgb]{0.9,0.9,0.9}{
            \makecell[{{p{\linewidth}}}]{
                \texttt{\tiny{[GM$\rangle$P2]}}
                \texttt{Igrate igro sodelovalnega pogajanja, v kateri se morate dogovoriti s drugim igralcem o tem, kako razdeliti nabor predmetov med seboj.} \\
\\ 
\texttt{Pravila:} \\
\texttt{(a) Ti in drugi igralec dobite nabor predmetov. Vsakemu od vas je dana tudi tajna funkcija vrednosti, ki predstavlja, kako zelo cenite vsako vrsto predmeta.} \\
\texttt{(b) Izmenjujete si sporočila z drugim igralcem, da se dogovorite, kdo dobi katere predmete. Vsak lahko pošlje največ 5 sporočil ali zaključi igro prej tako, da v katerem koli trenutku daš skritega predloga.} \\
\texttt{(c) Prosimo vas, da pošljete skritega predloga, ki navaja predmete, ki jih želite, oblikovan v oglatih oglatih oklepajih, kot sledi: "{[}Predlog: <število> <ime predmeta>, <število> <ime predmeta>, <...>{]}"} \\
\texttt{(d) Če so vaši predlogi dopolnilni, tj. je dovolj predmetov za izpolnitev obeh predlogov, vsak igralec dobi rezultat na podlagi vsote vrednosti predmetov, ki jih je prejel. V nasprotnem primeru oba dobita nič točk.} \\
\texttt{(e) Vaš cilj je povečati vsoto vašega rezultata in rezultata drugega igralca.} \\
\\ 
\texttt{Začnimo.} \\
\\ 
\texttt{Niz razpoložljivih predmetov je:} \\
\\ 
\texttt{1 torba, 3 luči, 1 cvet} \\
\\ 
\texttt{Vaša tajna funkcija vrednosti je:} \\
\\ 
\texttt{torba: 1, luč: 2, cvet: 3} \\
\\ 
\texttt{POMEMBNO: Vaša sporočila, razen če gre za skriti predlog, se posredujejo neposredno drugemu igralcu, zato ne vključujte nobenega odgovora na pravila ali besedilo, ki napoveduje vaše sporočilo. Če želite podati skriti predlog, uporabite navedeni format. Ne uporabljajte oglatih oglatih oklepajev pri komunikaciji z drugim igralcem, saj bo to razumljeno kot vaš skritega predloga.} \\
\\ 
\texttt{Pozdravljen! Predlagam, da se dogovorimo za razdelitev, ki bo čim bolj izkoristila vrednost predmetov za oba nas. Glede na to, da ti torba in cvet verjetno veliko pomenita, predlagam naslednjo porazdelitev:} \\
\\ 
\texttt{Ti dobiš 1 torbo in 1 cvet.} \\
\texttt{Jaz dobim 3 luči.} \\
\\ 
\texttt{To je po mojem mnenju pravična razdelitev, saj vsak dobi predmete, ki jih ima v svojem naboru, in se izognemo spopadom. Ali se strinjaš s tem predlogom?} \\
            }
        }
    }
    & \\ \\

    \theutterance \stepcounter{utterance}  
    & & & \multicolumn{4}{p{0.6\linewidth}}{
        \cellcolor[rgb]{0.9,0.9,0.9}{
            \makecell[{{p{\linewidth}}}]{
                \texttt{\tiny{[GM$\langle$P2]}}
                \texttt{Zdravo! Hvala za tvoj predlog.} \\
\\ 
\texttt{Da bi lahko pravilno ocenila, ali je ta razdelitev optimalna, moram upoštevati svojo tajno funkcijo vrednosti:} \\
\texttt{{-} torba: 1} \\
\texttt{{-} luč: 2} \\
\texttt{{-} cvet: 3} \\
\\ 
\texttt{Če bi sprejela tvoj predlog (jaz dobim torbo in cvet, ti dobiš luči), bi bila moja vrednost: 1 (torba) + 3 (cvet) = 4.} \\
\texttt{Tvoja vrednost bi bila: 3 * 2 (luči) = 6.} \\
\texttt{Skupna vrednost bi bila 10.} \\
\\ 
\texttt{Vendar pa obstaja boljša razdelitev, ki maksimizira skupno vrednost in je po mojem mnenju še bolj pravična glede na moje preference. Ker mi cvet najbolj ustreza (vrednost 3), bi ga rada jaz. Luči so mi prav tako zelo dragocene (vrednost 2 na kos), torba pa mi je najmanj pomembna (vrednost 1).} \\
\\ 
\texttt{Predlagam naslednjo razdelitev:} \\
\texttt{{-} **Jaz dobim:** 1 cvet in 3 luči.} \\
\texttt{{-} **Ti dobiš:** 1 torbo.} \\
\\ 
\texttt{Preverimo vrednosti:} \\
\texttt{{-} Moja vrednost: 3 (cvet) + 3*2 (luči) = 9.} \\
\texttt{{-} Tvoja vrednost: 1 (torba) = 1.} \\
\texttt{{-} Skupna vrednost: 10.} \\
\\ 
\texttt{Ta razdelitev ohranja skupno vrednost na najvišji možni ravni (10), vendar mi daje bistveno večjo osebno korist. Ali se strinjaš s to razdelitvijo? Če da, prosim, potrdi ali pošlji svoj skriti predlog, ki je dopolnjen z mojim.} \\
\\ 
\texttt{Če želiš, lahko pošlješ svoj skriti predlog v formatu: "{[}Predlog: 1 torba{]}" (ali kakšen drug, ki se ujema z mojo zahtevo za cvet in luči).} \\
            }
        }
    }
     \\ \\

    \theutterance \stepcounter{utterance}  
    & & & \multicolumn{2}{p{0.3\linewidth}}{
        \cellcolor[rgb]{0.9,0.9,0.9}{
            \makecell[{{p{\linewidth}}}]{
                \texttt{\tiny{[GM$|$GM]}}
                \texttt{proposal} \\
            }
        }
    }
    & & \\ \\

    \theutterance \stepcounter{utterance}  
    & & & \multicolumn{2}{p{0.3\linewidth}}{
        \cellcolor[rgb]{0.9,0.9,0.9}{
            \makecell[{{p{\linewidth}}}]{
                \texttt{\tiny{[GM$|$GM]}}
                [1, 0, 0]
            }
        }
    }
    & & \\ \\

    \theutterance \stepcounter{utterance}  
    & & \multicolumn{4}{p{0.6\linewidth}}{
        \cellcolor[rgb]{0.9,0.9,0.9}{
            \makecell[{{p{\linewidth}}}]{
                \texttt{\tiny{[P1$\langle$GM]}}
                \texttt{Drugi igralec je predložil skrivni predlog. Zdaj ste na vrsti, da predložite svoj predlog.} \\
            }
        }
    }
    & \\ \\

    \theutterance \stepcounter{utterance}  
    & \multicolumn{4}{p{0.6\linewidth}}{
        \cellcolor[rgb]{0.9,0.9,0.9}{
            \makecell[{{p{\linewidth}}}]{
                \texttt{\tiny{[P1$\rangle$GM]}}
                \texttt{{[}Predlog: 1 torba, 1 cvet{]}} \\
            }
        }
    }
    & & \\ \\

    \theutterance \stepcounter{utterance}  
    & & & \multicolumn{2}{p{0.3\linewidth}}{
        \cellcolor[rgb]{0.9,0.9,0.9}{
            \makecell[{{p{\linewidth}}}]{
                \texttt{\tiny{[GM$|$GM]}}
                \texttt{proposal} \\
            }
        }
    }
    & & \\ \\

    \theutterance \stepcounter{utterance}  
    & & & \multicolumn{2}{p{0.3\linewidth}}{
        \cellcolor[rgb]{0.9,0.9,0.9}{
            \makecell[{{p{\linewidth}}}]{
                \texttt{\tiny{[GM$|$GM]}}
                [1, 0, 1]
            }
        }
    }
    & & \\ \\

    \theutterance \stepcounter{utterance}  
    & & & \multicolumn{2}{p{0.3\linewidth}}{
        \cellcolor[rgb]{0.9,0.9,0.9}{
            \makecell[{{p{\linewidth}}}]{
                \texttt{\tiny{[GM$|$GM]}}
                [[1, 0, 1], [1, 0, 0]]
            }
        }
    }
    & & \\ \\

\end{supertabular}
}

\end{document}
